%% family_centroids.tex -- shared sub-recipe. Referenced by every analysis that
%% measures a gene against its family, via \dependson{sub:family-centroids}.

\begin{recipe}[label=sub:family-centroids]{Gene Family Centroids}

\recipepurpose{Reduce each gene family to a single expression vector, its
  centroid: the reference that later recipes measure divergence against.
  Distances to the centroid, shift vectors from it, and the outlier test built
  on those all inherit whichever set of family members you average here.}

\nomaintainer

\begin{requiredinputs}
  \item Normalized expression vectors, one gene per column
        (\code{(n\_axes, n\_genes)}), from the normalization sub-recipe.
  \item A gene-to-family mapping: for each gene, the index of its family, or
        \code{0} for a gene that belongs to none.
  \item Optional, and only for the orthologs-only centroid: which genes are
        the orthologs. This is \emph{your} data --- it comes from an orthology
        pipeline (OrthoFinder, synteny/Cactus, OMAmer), not from Tensor Omics.
\end{requiredinputs}

\begin{prerequisites}
  \dependson{sub:normalization}
  \dependson{sub:gene-family-detection}
\end{prerequisites}

\begin{workflow}
  \step{Number the families $1 \dots n$ and build the \code{gene\_to\_family}
        vector, one entry per column of the expression matrix. Use \code{0}
        for every gene that belongs to no family --- it is the sentinel the
        library expects, and such genes are skipped.}
  \step{Choose which members define the reference. \code{group\_centroid\_all}
        averages every gene of the family; \code{group\_centroid\_orthologs}
        averages only the genes marked in \code{ortholog\_set}. See the
        interpretation below --- this choice is the substance of the step,
        not a detail.}
  \step{Compute the centroids. The result is an
        \code{(n\_axes, n\_families)} matrix holding \emph{one centroid per
        column}, in the same layout and orientation as the expression vectors
        that went in.}
  \step{Count, per family, how many genes actually contributed. A family with
        no contributor is not reported by the library --- it silently receives
        a zero centroid, which is a valid point in the space and will quietly
        distort everything measured against it.}
\end{workflow}

\examplescript{python/examples/gene_family_centroids.py}{%
  Builds the gene-to-family mapping from an OrthoFinder-style family table,
  computes the centroids and writes them, reporting assignment coverage and
  any family left without a contributing gene. Run unchanged from the
  repository root, it computes all-genes centroids for the 458 bundled
  families over \file{material/normalization.tsv} and writes
  \file{results/family_centroids.tsv}; pass \optn{-{}-orthologs} with a list of
  gene ids to take the orthologs-only centroid instead.}

\begin{codefragment}[language=Python,caption={Fragments to edit: the mapping, and which members define the centroid}]
import numpy as np
from tensor_omics import (group_centroid_all,
                          group_centroid_orthologs)

# --- one gene per COLUMN, as normalization produced it -----
expr = np.asfortranarray(normalized)   # (n_axes, n_genes)

# --- family index per gene, 1-based; 0 = no family ---------
gene_to_family = np.array([1, 1, 2, 0, 2], dtype=np.int32)
n_families = 2

# --- (a) the family's actual mean, paralogs included -------
centroids = group_centroid_all(
    expr, gene_to_family, n_families)

# --- (b) a conserved reference: orthologs only -------------
ortholog_set = np.array([True, False, True, False, True])
centroids = group_centroid_orthologs(
    expr, gene_to_family, n_families, ortholog_set)
#   -> (n_axes, n_families), one centroid per column
\end{codefragment}

\begin{note}
  For a centroid over an arbitrary set of genes --- not a whole family ---
  call \code{mean\_vector(expr, gene\_indices)} directly. Its
  \code{gene\_indices} are 1-based column numbers, not a boolean mask.
\end{note}

\begin{parameters}
  \param{gene\_to\_family}{Which family each gene belongs to; entry $i$ is the
    family index of column $i$ of the expression matrix}{Values must be
    \code{1}\dots\code{n\_families}, or \code{0} for an unassigned gene. Any
    other value fails the call}
  \param{n\_families}{How many centroids are produced, i.e.\ the number of
    columns of the result}{Must be at least as large as the largest index in
    \code{gene\_to\_family}. Indices above it fail the call; families with no
    genes below it come back as zero centroids}
  \param{ortholog\_set}{Which genes contribute when the orthologs-only
    centroid is taken}{Required by \code{group\_centroid\_orthologs}, one
    boolean per gene. A family with no ortholog among its members gets a zero
    centroid}
\end{parameters}

\begin{expectedresults}
  The result is \code{(n\_axes, n\_families)}, read-only, and each column is
  the element-wise mean of the contributing genes --- so every centroid lies
  inside the range its family's members span on each axis. On the bundled
  data the script assigns 10\,234 of 85\,828 genes to 458 families (the rest
  belong to none), family sizes run from 6 to 143 genes, and no family is left
  empty. It also reports the 361 family members that have no row in the
  expression table; a much larger number there means the gene identifiers of
  the two files do not match.
\end{expectedresults}

\begin{pitfall}[Pitfall: a family with no members gets a zero centroid]
  A family that no gene contributes to --- because it has no members in the
  expression table, or none of them is an ortholog --- does not raise an
  error and is not flagged. It receives the zero vector, which is an ordinary
  point at the origin of the space. Every later distance to that centroid is
  then simply the gene's own norm, so the family looks maximally divergent for
  a purely bookkeeping reason. Count the contributors per family and check
  the count is never zero before going on.
\end{pitfall}

\begin{pitfall}[Pitfall: which members you average changes every later result]
  The two entry points are not variants of one number. On the bundled data
  the orthologs-only centroids shipped in
  \file{material/centroids_orthologs_only.tsv} sit a median 13.5\% of their
  own length away from the all-genes centroids of the same families, and up
  to 1.56 units away --- both computed from the same expression table, so the
  whole difference comes from the member set. Averaging \emph{all} genes folds
  the very paralogs you may later test for divergence into the reference they
  are tested against, which pulls the reference toward them and hides them.
  Use the orthologs-only centroid whenever the analysis asks whether a copy
  has diverged.
\end{pitfall}

\begin{pitfall}[Pitfall: marking unassigned genes with anything but zero]
  \code{0} is the sentinel for ``no family'', and genes carrying it are
  skipped. A negative index, or an index above \code{n\_families}, is not
  treated as ``unassigned'' but fails the call with
  \code{ToxInputError: invalid input arguments (argument 'gene\_to\_family')}.
  Map every gene you want excluded to \code{0} explicitly rather than leaving
  a placeholder behind.
\end{pitfall}

\begin{interpretation}
  A centroid is the family's consensus expression profile, and the distance
  of a member from it is that member's divergence from the family. What the
  consensus \emph{means} is decided by the member set. The orthologs-only
  centroid is a \emph{conserved} reference: it describes where the family's
  expression sat before the copies you are testing diverged, which is what
  makes a large distance interpretable as divergence of that copy. The
  all-genes centroid is the family's actual centre of mass, including every
  divergent copy --- the right reference when you want to describe the family
  as it is, and the wrong one when you want to detect a member that moved.

  Read a centroid alongside its contributor count. A centroid over two genes
  is a midpoint, not a consensus, and distances measured against it carry the
  noise of those two genes.
\end{interpretation}

\begin{references}
  \reference{Emms \& Kelly (2019). OrthoFinder: phylogenetic orthology
    inference for comparative genomics. \emph{Genome Biology} 20:238.}
\end{references}

\end{recipe}
